%% ch02_ibr_foundations.tex
%% Chapter: IBR Characterisation and Theoretical Foundations
%% Merged from: ch02_mathematical_foundations.tex, ch03_ibr_characterisation.tex
%% ============================================================
\chapter{IBR Characterisation and Theoretical Foundations}
\label{ch:foundations}
\label{ch:ibr_char}

\section{Introduction}
\label{sec:ch2_intro}

Every adaptive protection algorithm in this thesis rests on a common
mathematical structure: an \emph{inverse estimation problem}, in which
a physical model of the protected device is fitted to real-time measurements,
and the identified parameters are used to update relay thresholds. This
chapter establishes the theoretical foundations of that structure:
the general nonlinear least-squares formulation (\S\ref{sec:ch2_lsq}),
the Levenberg--Marquardt algorithm used as the solver
(\S\ref{sec:ch2_lm}), the composite confidence score that gates adaptation
decisions (\S\ref{sec:ch2_confidence}), and the bounded update law that
guarantees protection security (\S\ref{sec:ch2_update}).

These foundations are drawn upon directly by the generator OC
(Chapter~\ref{ch:generator_oc}) and transformer differential
(Chapter~\ref{ch:transformer_diff}) algorithms, and are reflected in
the EKF formulation of the digital twin (Chapter~\ref{ch:digital_twin}).

\section{Inverse Estimation Formulation}
\label{sec:ch2_lsq}

\subsection{General Structure}

Let $\mathbf{x}(t; \bm{\theta}) \in \mathbb{R}^n$ be the state vector
of a protected device (fault current, voltage, or impedance trajectory)
parameterised by an unknown vector $\bm{\theta} \in \mathbb{R}^p$.
Let $\mathbf{y}_k \in \mathbb{R}^n$ be a sequence of $K$ noisy measurements
at times $t_k$. The \emph{inverse estimation problem} is:

\begin{equation}
  \hat{\bm{\theta}} = \argmin_{\bm{\theta} \in \Theta}
    \sum_{k=1}^{K} w_k \left\|\mathbf{y}_k - \mathbf{x}(t_k;\bm{\theta})\right\|^2
  \label{eq:ch2_lsq}
\end{equation}

where $\Theta \subset \mathbb{R}^p$ is the feasible set defined by physical
bounds, $w_k > 0$ are sample weights (typically $w_k = 1$ for uniform
confidence within the measurement window), and $\|\cdot\|$ is the Euclidean
norm.

Writing $\mathbf{r}(\bm{\theta}) = \mathbf{y} - \mathbf{x}(\bm{\theta})$
as the residual vector (stacking all $K$ time steps), the problem becomes:
\begin{equation}
  \hat{\bm{\theta}} = \argmin_{\bm{\theta} \in \Theta}
    \frac{1}{2}\|\mathbf{r}(\bm{\theta})\|^2
  \label{eq:ch2_rss}
\end{equation}

When $\mathbf{x}(t;\bm{\theta})$ is nonlinear in $\bm{\theta}$ (as is the
case for all physical models in this thesis), Eq.~\eqref{eq:ch2_rss} is a
\emph{bounded nonlinear least-squares} problem.

\subsection{Jacobian and Sensitivity Analysis}

The first-order necessary condition $\nabla_{\bm{\theta}}(\frac{1}{2}\|\mathbf{r}\|^2) = 0$
gives the \emph{normal equations}:
\begin{equation}
  J(\bm{\theta})^\top \mathbf{r}(\bm{\theta}) = \mathbf{0}
  \label{eq:ch2_normal}
\end{equation}
where $J \in \mathbb{R}^{(nK) \times p}$ is the Jacobian:
\begin{equation}
  J_{ij} = \frac{\partial x_i}{\partial \theta_j}
  \label{eq:ch2_jacobian}
\end{equation}

The condition number $\kappa(J)$ measures the sensitivity of the solution
to measurement noise. When columns of $J$ are nearly linearly dependent
(ill-conditioning), small measurement perturbations cause large parameter
errors. Column-normalising $J$ before evaluating $\kappa$ removes
dimension effects:
\begin{equation}
  \kappa_n(J) = \kappa\left(\mathrm{diag}(\|J_{\cdot,1}\|,\ldots,\|J_{\cdot,p}\|)^{-1} J\right)
  \label{eq:ch2_kappan}
\end{equation}

\section{The Levenberg--Marquardt Algorithm}
\label{sec:ch2_lm}

\subsection{Derivation}

The Levenberg--Marquardt (LM) algorithm~\citep{Levenberg1944,Marquardt1963}
interpolates between steepest-descent (far from minimum) and
Gauss--Newton (near minimum) by adding a trust-region regulariser:

\begin{equation}
  \bm{\delta}^{(i)} = -(J^\top J + \mu^{(i)} I)^{-1} J^\top \mathbf{r}
  \label{eq:ch2_lm_step}
\end{equation}

\begin{equation}
  \bm{\theta}^{(i+1)} = \Pi_\Theta\left(\bm{\theta}^{(i)} + \bm{\delta}^{(i)}\right)
  \label{eq:ch2_lm_update}
\end{equation}

where $\Pi_\Theta(\cdot)$ projects onto the feasible set $\Theta$ (element-wise
clipping to physical bounds), and $\mu^{(i)} > 0$ is the damping parameter.

\subsection{Damping Parameter Adaptation}

The damping parameter is adapted by the gain ratio:
\begin{equation}
  \rho^{(i)} = \frac{\|\mathbf{r}(\bm{\theta}^{(i)})\|^2
                   - \|\mathbf{r}(\bm{\theta}^{(i+1)})\|^2}
                    {L(\bm{\theta}^{(i)}) - L(\bm{\theta}^{(i+1)})}
  \label{eq:ch2_gain_ratio}
\end{equation}
where $L(\bm{\theta})$ is the linearised model prediction:
$L(\bm{\theta}^{(i)}) - L(\bm{\theta}^{(i+1)}) = \bm{\delta}^\top(2J^\top \mathbf{r} - J^\top J\bm{\delta})$.

The update rule is:
\begin{equation}
  \mu^{(i+1)} = \begin{cases}
    \mu^{(i)} / 3  & \text{if } \rho^{(i)} > 0.75 \text{ (good step: reduce damping)}\\
    \mu^{(i)}      & \text{if } 0.25 \leq \rho^{(i)} \leq 0.75 \\
    \mu^{(i)} \times 10 & \text{if } \rho^{(i)} < 0.25 \text{ (poor step: increase damping)}
  \end{cases}
  \label{eq:ch2_mu_adapt}
\end{equation}

\subsection{Two-Pass Strategy}

For physical models where parameters separate into ``fast'' (identifiable
early in the measurement window) and ``slow'' (identifiable only at the tail
of the window) components, a two-pass approach improves convergence:

\begin{enumerate}
\item \textbf{Pass 1 (transient identification):} Fit the model over the
      early window $[t_0, t_0 + T_1]$ with the slow parameters held fixed at
      nominal values. Estimate the transient parameters $\bm{\theta}_\mathrm{fast}$.
\item \textbf{Pass 2 (steady-state identification):} Fit over the full
      window $[t_0, t_0 + T_2]$ with $\bm{\theta}_\mathrm{fast}$ fixed at
      Pass~1 values. Estimate the slow parameters $\bm{\theta}_\mathrm{slow}$.
\end{enumerate}

This strategy reduces the effective Jacobian condition number by decoupling
the transient and steady-state parameter subspaces, and is used in
the SAMBP-SyncOC algorithm (Chapter~\ref{ch:generator_oc}).

\begin{algorithm}[htbp]
\caption{Levenberg--Marquardt Inverse Estimator}
\label{alg:ch2_lm}
\begin{algorithmic}[1]
\Require Measurement window $\mathbf{y}$, model $\mathbf{x}(\cdot;\bm{\theta})$,
         initial guess $\bm{\theta}^{(0)}$, bounds $\Theta$,
         $\mu_0$, $\mu_\mathrm{max}$, $N_\mathrm{iter}$, $\varepsilon_r$
\Ensure  $\hat{\bm{\theta}}$, residual $R^2$, Jacobian $J$
\State $i \gets 0$,\; $\mu \gets \mu_0$
\While{$i < N_\mathrm{iter}$ and $\|\mathbf{r}^{(i)}\| > \varepsilon_r$}
  \State Compute $\mathbf{r}^{(i)} = \mathbf{y} - \mathbf{x}(\bm{\theta}^{(i)})$
  \State Compute Jacobian $J^{(i)}$ via finite differences
  \State Solve $(J^\top J + \mu I)\bm{\delta} = -J^\top \mathbf{r}^{(i)}$
  \State $\bm{\theta}^{\mathrm{cand}} \gets \Pi_\Theta(\bm{\theta}^{(i)} + \bm{\delta})$
  \State Compute gain ratio $\rho$ \Comment{Eq.~\eqref{eq:ch2_gain_ratio}}
  \If{$\rho > 0$}
    \State $\bm{\theta}^{(i+1)} \gets \bm{\theta}^{\mathrm{cand}}$
  \Else
    \State $\bm{\theta}^{(i+1)} \gets \bm{\theta}^{(i)}$ \Comment{reject step}
  \EndIf
  \State Update $\mu$ per Eq.~\eqref{eq:ch2_mu_adapt};\; $i \gets i+1$
\EndWhile
\State $\hat{\bm{\theta}} \gets \bm{\theta}^{(i)}$
\State Compute $R^2 = 1 - \|\mathbf{r}^{(i)}\|^2 / \|\mathbf{y} - \bar{\mathbf{y}}\|^2$
\end{algorithmic}
\end{algorithm}

\subsection{Convergence Properties}

Under standard regularity conditions~\citep{Nocedal2006}, the LM algorithm
satisfies:
\begin{equation}
  \|\bm{\theta}^{(i+1)} - \hat{\bm{\theta}}\| \leq C \|\bm{\theta}^{(i)} - \hat{\bm{\theta}}\|^{1+q}
  \label{eq:ch2_lm_convergence}
\end{equation}
with $q \in [1, 2)$ (superlinear convergence near the minimum, when $\mu \to 0$).
For the protection application, convergence to $R^2 > 0.999$ within 15--20
iterations has been demonstrated across all study cases (Chapter~\ref{ch:generator_oc}).

\section{Composite Confidence Score}
\label{sec:ch2_confidence}

After solving the inverse estimation problem, it is essential to know
\emph{how much to trust the result} before allowing it to update relay
settings. A wrong estimate with high apparent $R^2$ can arise from
overfitting, poor identifiability, or measurement noise. A four-component
composite confidence score $\gamma \in [0,1]$ is defined as:

\begin{equation}
  \gamma = w_1 \gamma_R + w_2 \gamma_J + w_3 \gamma_B + w_4 \gamma_P
  \label{eq:ch2_gamma}
\end{equation}

with $w_1 + w_2 + w_3 + w_4 = 1$ (default: $w = [0.35, 0.30, 0.20, 0.15]$).

\paragraph{Component $\gamma_R$ --- Normalised Residual Quality.}
\begin{equation}
  \gamma_R = \frac{1}{1 + \alpha_R (1-R^2)}
  \label{eq:ch2_gamma_r}
\end{equation}
where $\alpha_R = 100$ sharpens the score at $R^2 \approx 1$.
$\gamma_R \to 1$ when $R^2 \to 1$ (near-perfect fit).

\paragraph{Component $\gamma_J$ --- Jacobian Condition Number.}
\begin{equation}
  \gamma_J = \frac{1}{1 + \alpha_J \log_{10}(\kappa_n(J))}
  \label{eq:ch2_gamma_j}
\end{equation}
with $\alpha_J = 0.1$. $\gamma_J \to 1$ when $\kappa_n(J) \to 1$
(well-conditioned, identifiable problem); $\gamma_J \to 0$ as
$\kappa_n(J) \to \infty$ (unidentifiable).

\paragraph{Component $\gamma_B$ --- Bounds-Hit Fraction.}
\begin{equation}
  \gamma_B = 1 - \frac{N_\mathrm{bounds}}{p}
  \label{eq:ch2_gamma_b}
\end{equation}
where $N_\mathrm{bounds}$ is the number of parameters that converged to
a boundary of $\Theta$. Parameters hitting bounds indicate that the
optimum is physically infeasible or the model is misidentified.

\paragraph{Component $\gamma_P$ --- Physics Plausibility.}
$\gamma_P \in \{0, 0.5, 1.0\}$ is a rule-based score on the
physical plausibility of $\hat{\bm{\theta}}$. Each estimated parameter
is checked against physics-based secondary bounds:
\begin{equation}
  \gamma_P = 1 - \frac{N_\mathrm{implausible}}{p}
  \label{eq:ch2_gamma_p}
\end{equation}
where $N_\mathrm{implausible}$ is the number of parameters outside
physically expected ranges (e.g., $\hat{X}_d'' > \hat{X}_d'$,
$\hat{\tau}_{ac} < 0$).

\paragraph{Adaptation Gate.}
Adaptation of relay settings is authorised only when:
\begin{equation}
  \gamma \geq \gamma_\mathrm{th}
  \label{eq:ch2_gate}
\end{equation}
with default $\gamma_\mathrm{th} = 0.70$.
This threshold was chosen to balance responsiveness (adapt quickly to
real changes) against security (reject spurious estimates).

\section{Bounded Adaptive Update Law}
\label{sec:ch2_update}

\subsection{Motivation}

If the inverse estimator produces an erroneous $\hat{\bm{\theta}}$ that
passes the confidence gate (false positive), the relay settings must not
move to an unsafe operating point. The bounded update law imposes three
levels of safety:

\begin{enumerate}
\item \textbf{Rate limiting:} Settings change at most $\delta_\mathrm{max}$ per
      update cycle, preventing step changes.
\item \textbf{Hard bounds:} Settings are clamped to $[\theta_\mathrm{min},
      \theta_\mathrm{max}]$ from commissioning specifications.
\item \textbf{Security floor:} Settings may never drop below the minimum
      value required to detect the worst-case bolted fault with a 20\%
      margin.
\end{enumerate}

\subsection{Update Equation}

Let $\xi_k$ be the $k$-th relay setting (e.g., pickup $I_p$ or
time-multiplier $\TMS$), and $\hat{\xi}_k(\hat{\bm{\theta}})$ the
value recommended by the identified model. The bounded update is:

\begin{equation}
  \xi_k^{(n+1)} = \mathrm{clip}\!\left(
    \xi_k^{(n)} + \delta_\mathrm{max}\,
    \mathrm{sgn}(\hat{\xi}_k - \xi_k^{(n)})
    \cdot \min\left(1,\,
      \frac{|\hat{\xi}_k - \xi_k^{(n)}|}{\delta_\mathrm{max}}
    \right),\;
    \xi_\mathrm{min},\; \xi_\mathrm{max}
  \right)
  \label{eq:ch2_update}
\end{equation}

For the IEC Very-Inverse relay:
\begin{align}
  I_p^{(n+1)} &= \mathrm{clip}\!\left(I_p^{(n)} + \Delta I_p,\,
    1.05 I_\mathrm{load}^{\max},\, 0.80 I_\mathrm{fault}^{\min}\right)
  \label{eq:ch2_ip_update} \\
  \TMS^{(n+1)} &= \mathrm{clip}\!\left(\TMS^{(n)} + \Delta \TMS,\,
    \TMS_\mathrm{min},\, \TMS_\mathrm{max}\right)
  \label{eq:ch2_tms_update}
\end{align}

\subsection{Security Proof}

\begin{proposition}[Bounded update preserves trip security]
\label{prop:ch2_security}
If $\xi_k^{(0)}$ is in $[\xi_\mathrm{min}, \xi_\mathrm{max}]$ and all
$\hat{\xi}_k$ are in $\mathbb{R}$, then for all $n \geq 0$:
$\xi_k^{(n)} \in [\xi_\mathrm{min}, \xi_\mathrm{max}]$.
\end{proposition}

\begin{proof}
The update Eq.~\eqref{eq:ch2_update} applies $\mathrm{clip}(\cdot,
\xi_\mathrm{min}, \xi_\mathrm{max})$ at every step. Since $\mathrm{clip}$
is idempotent and bounds-preserving, the invariant holds by induction.
\end{proof}

\section{Applicability Map}
\label{sec:ch2_map}

Table~\ref{tab:ch2_applicability} summarises how the inverse estimation
framework specialises to each protection element in this thesis.
The mathematical structure (Eq.~\eqref{eq:ch2_lsq}--\eqref{eq:ch2_update})
is identical; only the physical model $\mathbf{x}(t;\bm{\theta})$ and
the parameter vector $\bm{\theta}$ change.

\begin{table}[htbp]
\caption{Inverse estimation framework specialisation by protection element}
\label{tab:ch2_applicability}
\small
\begin{tabularx}{\textwidth}{@{}lp{2.5cm}p{3.5cm}p{2.5cm}p{2cm}@{}}
\toprule
Element & Model & Parameters $\bm{\theta}$ & Adapted setting & Status \\
\midrule
87L line diff & IBR differential current & $k_\mathrm{ibr}$, $\phi_B$ & Bias slope $k_m^*$ & [P] \\
Gen. OC (SG) & 3-component fault current & $\{I_{ss},I_{sub},\tau_{ac},\tau_{dc},\phi_a\}$ & $I_p$, $\TMS$ & [P] \\
Gen. OC (DFIG) & Piecewise crowbar model & $\{I_{0+},\tau_r,I_{crowd}\}$ & $I_p$ \\
87T transformer & Inrush vs.\ fault model & $\{i_2/i_1,\phi_\mathrm{hm}\}$ & SPRT threshold \\
EKF digital twin & IBR Thevenin model & $\kibr$, $\Zline$, $\alpha$ & 87L bias $k_m^*$ & [P] \\
\bottomrule
\end{tabularx}
\end{table}

\section{Relationship to Kalman Filtering}
\label{sec:ch2_ekf_relation}

The EKF used in the digital twin (Chapter~\ref{ch:digital_twin}) is
formally a \emph{recursive} solution to the inverse estimation problem
with Gaussian process and measurement noise. Writing the state model as:
\begin{align}
  \mathbf{x}_k &= f(\mathbf{x}_{k-1}) + \mathbf{w}_k, \quad
    \mathbf{w}_k \sim \mathcal{N}(\mathbf{0}, Q) \\
  \mathbf{y}_k &= h(\mathbf{x}_k) + \mathbf{v}_k, \quad
    \mathbf{v}_k \sim \mathcal{N}(\mathbf{0}, R)
\end{align}
the EKF update step minimises the same objective as
Eq.~\eqref{eq:ch2_lsq} over a sliding window, with the prediction
step providing the Tikhonov regularisation that plays the role of the
LM damping term $\mu$. This formal connection means the confidence
score framework of \S\ref{sec:ch2_confidence} applies directly to
evaluate the quality of EKF state estimates.



%% ── Part 2: IBR Fault Current Characterisation and Protection Challenges ──────────────────────────────────

\section{Introduction}
\label{sec:ch3_intro}

An accurate mathematical model of each generation source is the prerequisite
for designing protection algorithms that respect physical constraints imposed
by power electronic interfaces. This chapter develops the models used throughout
the thesis, quantifies how IBR fault signatures degrade legacy relay performance,
and motivates the adaptive solutions proposed in Chapters~\ref{ch:line_diff}--\ref{ch:centralised}.

Four source types are treated:
\begin{enumerate}
\item \textbf{Synchronous Generator (SG):} The classical three-component
      subtransient fault current with DC offset (Section~\ref{sec:ch3_sg}).
\item \textbf{Grid-Following Inverter (GFL):} Current-controlled VSC with
      fast current limiter (Section~\ref{sec:ch3_gfl}).
\item \textbf{Grid-Forming Inverter (GFM):} Voltage-controlled VSC emulating
      synchronous machine inertia (Section~\ref{sec:ch3_gfm}).
\item \textbf{DFIG Type-3 and PMSG Type-4 Wind:} Models planned for future
      extension of the SAMBP framework (Section~\ref{sec:ch3_dfig}).
\end{enumerate}

\section{Synchronous Generator Fault Current Model}
\label{sec:ch3_sg}

\subsection{Three-Component Fault Current}

The Thevenin equivalent at the SG terminals for a three-phase fault at $t=0^+$ is:

\begin{equation}
  i_a^\mathrm{SG}(t) =
    \frac{\sqrt{2} E_0}{X_d''} e^{-t/T_d''}
    + \left(\frac{1}{X_d'} - \frac{1}{X_d''}\right)\sqrt{2} E_0\, e^{-t/T_d'}
    + \left(\frac{1}{X_d} - \frac{1}{X_d'}\right)\sqrt{2} E_0
    + I_\mathrm{DC}\,e^{-t/T_a}
  \label{eq:ch3_sg_fault}
\end{equation}

where $X_d''$, $X_d'$, $X_d$ are sub-transient, transient, and synchronous
direct-axis reactances; $T_d''$, $T_d'$ are the corresponding time constants;
$T_a = X''_d / (\omega_0 R_a)$ is the DC offset time constant; and
$I_\mathrm{DC} = \sqrt{2} E_0 / X_d'' \cdot \cos(\phi_0)$ depends on the
pre-fault power angle $\phi_0$. For a typical large SG, the initial peak
fault current is 6--10~pu, decaying over hundreds of milliseconds.

The \textbf{reduced-order model} used in Chapter~\ref{ch:generator_oc}
for the inverse estimation problem collapses Eq.~\eqref{eq:ch3_sg_fault}
to six identifiable parameters:
\begin{equation}
  i_a(t) = I_{ss} + I_{sub}\,e^{-(t-t_0)/\tau_{ac}}\cos(\omega_0 t + \phi_a)
           + I_{DC}(t_0)\,e^{-(t-t_0)/\tau_{dc}}
  \label{eq:ch3_sg_reduced}
\end{equation}
where $\bm{\theta} = \{t_0, I_{ss}, I_{sub}, \tau_{ac}, \tau_{dc}, \phi_a\}$
and the DC amplitude $I_{DC}(t_0) = -I_{sub}\cos\phi_a - I_{ss}$
is fixed by the continuity condition at $t = t_0^+$, eliminating one
free parameter from the otherwise seven-parameter formulation.

\subsection{Sequence Characteristics}

Under unbalanced faults, the SG exhibits $Z_1 \neq Z_2 \neq Z_0$:
\begin{equation}
  Z_1^\mathrm{SG} = \jj X_d'', \quad
  Z_2^\mathrm{SG} = \jj \frac{X_d'' + X_q''}{2}, \quad
  Z_0^\mathrm{SG} = \jj X_0 + 3Z_n
  \label{eq:ch3_sg_seq_z}
\end{equation}
The non-zero $Z_2$ and $Z_0$ drive negative-sequence and zero-sequence
currents during SLG faults, which are the physical basis for the 87LN
and 87LN0 elements of Chapter~\ref{ch:line_diff}.

\section{Grid-Following Inverter Model}
\label{sec:ch3_gfl}

\subsection{Inner Current Loop Dynamics}

In the rotating $dq$ reference frame, the converter filter equations are:
\begin{align}
  L_f \frac{d\id}{dt} &= -R_f \id + \omega_0 L_f \iq + v_d - e_d \\
  L_f \frac{d\iq}{dt} &= -R_f \iq - \omega_0 L_f \id + v_q - e_q
\end{align}

The closed-loop bandwidth $\omega_\mathrm{bw} \approx 2\pi \times 500$~rad/s
gives a current step response settling time of $\approx 2$~ms.

\subsection{Current Limiter and Fault Current}

The IGBT thermal protection clips $|i_{dq}| \leq I_\mathrm{lim}$:
\begin{equation}
  \phasor{I}_\mathrm{GFL} = I_\mathrm{lim}\angle(\theta_\mathrm{PLL} + \phi_\mathrm{cmd}),
  \quad I_\mathrm{lim} \in [1.1, 1.5]\,I_\mathrm{rated}
\end{equation}

The \emph{IBR current ratio} is the key protection parameter used throughout
this thesis:
\begin{equation}
  \kibr \triangleq \frac{|\phasor{I}_\mathrm{GFL}|}{I_\mathrm{rated}}
  \in [0.03, 0.55]
  \label{eq:ch3_kibr_def}
\end{equation}

\subsection{Sequence Characteristics of GFL}

Under IEEE~2800 negative-sequence suppression (default):
\begin{equation}
  \phasor{I}_2^\mathrm{GFL} = 0, \qquad
  \phasor{I}_0^\mathrm{GFL} = 0 \quad \text{(3-wire converter)}
  \label{eq:ch3_gfl_seq}
\end{equation}

This has critical implications: the 87LN element (negative-sequence
differential) does \emph{not} depend on the IBR fault current contribution,
making it IBR-resilient by design (Chapter~\ref{ch:line_diff}).

\section{Grid-Forming Inverter Model}
\label{sec:ch3_gfm}

\subsection{Virtual Synchronous Machine (VSM)}

The GFM emulates the swing equation:
\begin{equation}
  M\ddot{\theta} + D\dot{\theta} = P^\ast - P_e
\end{equation}
with virtual impedance $Z_v = R_v + \jj\omega_0 L_v$ giving fault current:
\begin{equation}
  I_\mathrm{sc}^\mathrm{GFM} = \frac{|V_0|}{|Z_v|} \approx \frac{|V_0|}{\omega_0 L_v}
  \in [1.5, 2.5]\,I_\mathrm{rated}
\end{equation}

GFM fault current is higher than GFL but lower than SG. The VSM model
is validated against 6th-order Park $dq$ simulations for $t \leq 150$~ms.
For $t > 150$~ms, AVR dynamics become significant and the reduced model
diverges by up to 5\%.

\begin{table}[htbp]
\caption{Fault current characteristics: SG vs.\ GFL vs.\ GFM}
\label{tab:ch3_fault_comparison}
\begin{tabularx}{\textwidth}{@{}lp{3cm}p{3cm}p{3cm}X@{}}
\toprule
Parameter & SG & GFL & GFM (VSM) & Relay impact \\
\midrule
Magnitude (pu) & 6--10 & 1.1--1.5 & 1.5--2.5
  & OC pickup; 87L bias slope \\
Decay & Exponential ($T_d'' \!\approx\! 30$~ms) & Controlled constant
  & Controlled constant
  & TDCS detection window \\
$I_2$ injection & $\propto V_2$ (physical) & Zero (suppressed)
  & Zero (suppressed)
  & 87LN resilience to IBR \\
$I_0$ injection & $\propto V_0$ (grounding) & Zero (3-wire)
  & Zero (3-wire)
  & 87LN0; 21G elements \\
Phase angle & Fixed (rotor) & Commanded $\pm 90$
  & Virtual rotor
  & Directional element \\
Response time & 5--50~ms & 2--5~ms & 5--20~ms
  & TDCS time window \\
\bottomrule
\end{tabularx}
\end{table}

\section{DFIG Type-3 and PMSG Type-4 Models (Extended Framework)}
\label{sec:ch3_dfig}

The DFIG and PMSG wind generators represent the dominant wind turbine
technologies globally, yet their fault current models are substantially
more complex than the GFL/GFM VSC models. This section presents the
theoretical models that will underpin the DFIG generator OC extension

\subsection{DFIG Type-3 Fault Current: Crowbar-Switched Piecewise Model}

The DFIG Type-3 wind turbine has a partially-rated rotor-side converter
(RSC) that injects current through the wound rotor. During a voltage dip,
the DC bus overvoltage protection activates the \emph{crowbar}---a
resistive shorting of the rotor circuit---within 2--5~ms. The fault
current exhibits two distinct phases:

\textbf{Phase I (pre-crowbar, $0 < t < t_{cb}$):}
\begin{equation}
  i_s(t) = I_{s0}\,e^{-t/T_s'}\cos(\omega_0 t + \phi_0)
          + I_{r,\mathrm{RSC}}\,e^{-t/T_r}
  \label{eq:ch3_dfig_phase1}
\end{equation}
where $T_s' = L_s'/R_s$ and the RSC current $I_{r,\mathrm{RSC}}$ is
limited by the converter current controller.

\textbf{Phase II (post-crowbar, $t \geq t_{cb}$):}
\begin{equation}
  i_s(t) = I_{s,cb}\,e^{-(t-t_{cb})/T_{s,cb}}\cos(\omega_0 t + \phi_{cb})
          + I_{s,dc}\,e^{-(t-t_{cb})/T_{dc}}
  \label{eq:ch3_dfig_phase2}
\end{equation}
with $T_{s,cb} = L_s / (R_s + R_{cb})$ where $R_{cb}$ is the crowbar
resistance. The crowbar activation introduces a step discontinuity at
$t = t_{cb}$ that is not captured by the SG reduced-order model.

The \textbf{DFIG inverse estimation problem} parameterised as:
\begin{equation}
  \bm{\theta}_\mathrm{DFIG} = \{t_0, t_{cb}, I_{s0}, I_{r,\mathrm{RSC}},
    T_s', T_r, I_{s,cb}, T_{s,cb}\}
  \label{eq:ch3_dfig_params}
\end{equation}
is an 8-parameter two-phase LM problem with the phase transition time
$t_{cb}$ as an additional free parameter. The two-pass LM strategy of
Chapter~\ref{ch:foundations} generalises directly: Pass~1 identifies
$\{I_{s0}, T_s', T_r\}$ over Phase~I, and Pass~2 identifies
$\{I_{s,cb}, T_{s,cb}\}$ over Phase~II with $t_{cb}$ from Pass~1.


\subsection{PMSG Type-4 Model}

The PMSG Type-4 connects to the grid through a full-converter back-to-back
VSC. The grid-side converter (GSC) is functionally identical to the GFL
model (Section~\ref{sec:ch3_gfl}): the fault current is
$|\phasor{I}_\mathrm{PMSG}| = I_\mathrm{lim}^{GSC} \in [1.1, 1.5]$~pu.
The protection implications are therefore identical to GFL, and no new
relay algorithm development is required for PMSG protection beyond the
framework of Chapter~\ref{ch:line_diff}.

\section{Protection Challenges Quantified}
\label{sec:ch3_challenges}

\subsection{Overcurrent Relay: Pickup Failure at Low $\kibr$}

The minimum detectable fault current with a standard ITOC relay is
$I_\mathrm{pickup} = k_s I_\mathrm{load,max}$ with $k_s = 1.3$.
For a pure-IBR feeder at maximum load ($I_\mathrm{load} \approx I_\mathrm{rated}$)
and minimum fault ($I_\mathrm{fault}^\mathrm{min} = 1.1\,I_\mathrm{rated}$):

\begin{equation}
  I_\mathrm{pickup} = 1.3\,I_\mathrm{rated} > I_\mathrm{fault}^\mathrm{min} = 1.1\,I_\mathrm{rated}
\end{equation}

The ITOC relay is \textbf{incapable of detecting faults} on IBR-only feeders
under maximum load. The adaptive OC framework of Chapter~\ref{ch:generator_oc}
addresses this for generator protection by tracking the generator model
rather than using a fixed pickup.

\subsection{87L Differential Relay: Bias Slope Conflict}

The classical 87L relay operates on:
$I_\mathrm{diff} > k_m I_\mathrm{bias} + I_0$.
For an internal fault with IBR at end~B:
\begin{equation}
  I_\mathrm{diff} = \sqrt{I_A^2 + I_B^2 + 2 I_A I_B \cos(\phi_A - \phi_B)}
\end{equation}
When $\kibr$ is low, $I_B \ll I_A$, so $I_\mathrm{diff} \approx I_A$ and the
relay is sensitive. When $\kibr$ is high (IBR at rated during ride-through)
but the phase angle $\phi_A - \phi_B \approx 90$ (reactive support mode),
$I_\mathrm{diff}$ drops to $\sqrt{I_A^2 + I_B^2}$, potentially below
$k_m I_\mathrm{bias} = k_m (I_A + I_B)/2$. The adaptive threshold of
Chapter~\ref{ch:line_diff} resolves this conflict.

\subsection{Distance Relay: Reach Limitation from IBR Reactive Injection}

The apparent impedance seen by a distance relay at end~A with IBR at end~B:
\begin{equation}
  \Zapp = \alpha \Zline + \frac{I_A + I_B}{I_A}\, Z_F
  \label{eq:ch3_zapp}
\end{equation}

The IBR injects reactive current $I_B = I_\mathrm{lim}\angle 90$,
shifting $\Zapp$ in the reactive direction on the R-X plane. For a
purely resistive fault ($Z_F = R_F$), the shift is:
\begin{equation}
  \Delta Z^\mathrm{react} = \frac{I_B^\mathrm{react}}{I_A}\, R_F
  \label{eq:ch3_reach_shift}
\end{equation}
This causes the relay to see $|\Zapp| > |\alpha \Zline|$ and underreach
(miss the fault). The quadrilateral correction of
Chapter~\ref{ch:distance} subtracts this shift using the DT estimate of $\kibr$.

\subsection{Transformer Differential: Inrush Mis-Discrimination at Low IBR}

The classical 87T relay uses second-harmonic blocking to restrain during
transformer energisation inrush. The threshold is set at:
$I_{2f}/I_{1f} > 0.15$ (15\% second harmonic ratio).
When IBRs supply the energisation current, the high-frequency harmonics
generated by the switching converters can suppress the inrush signature
below the 15\% threshold, causing the 87T to \emph{operate} (false trip)
during energisation. Conversely, internal turn-to-turn faults may generate
a small harmonic content that incorrectly blocks the 87T.
The SPRT solution of Chapter~\ref{ch:transformer_diff} replaces the fixed
harmonic threshold with a sequential likelihood ratio test that adapts to
the observed harmonic envelope.

\subsection{Summary of Failure Modes}

\begin{table}[htbp]
\caption{Legacy relay failure modes in IBR-dominated networks}
\label{tab:ch3_failures}
\small
\begin{tabularx}{\textwidth}{@{}lp{3cm}p{3.5cm}X@{}}
\toprule
Failure Mode & Element & Root Cause & SAMBP Solution \\
\midrule
F1 & ITOC pickup failure & $I_\mathrm{fault} < I_\mathrm{pickup}$ & Adaptive OC (Ch.~\ref{ch:generator_oc}) \\
F2 & 87L bias conflict & Fixed $k_m$, variable $\kibr$ & Adaptive 87L (Ch.~\ref{ch:line_diff}) \\
F3 & 87L SLG miss & IBR suppresses $I_2$ & 87LN + 87LN0 (Ch.~\ref{ch:line_diff}) \\
F4 & 87T inrush mis-discrim. & IBR distorts harmonic ratio & SPRT 87T (Ch.~\ref{ch:transformer_diff}) \\
F5 & Distance underreach & IBR reactive injection & Quadrilateral 21 (Ch.~\ref{ch:distance}) \\
F6 & Islanding blind spot & Mode-switch timing & Microgrid prot.\ (Ch.~\ref{ch:microgrid}) \\
F7 & PINN physics violation & No constraint in loss & PINN + KVL/KCL (Ch.~\ref{ch:ml}) \\
F8 & No closed-loop adaptation & Manual settings process & WAPC protocol (Ch.~\ref{ch:centralised}) \\
\bottomrule
\end{tabularx}
\end{table}
